% applications/discussion-DE.tex
\chapter{Diskussion}
\label{chap:applications-discussion}

Diese Diskussion bewertet Umfang, Stärken und Grenzen des Anwendungsteils
als Konsequenzebene genehmigter AIM-Begriffe.

\section{Einordnung}

Anwendungen ersetzen weder Grundlagen noch Zustand, Laufzeit oder
Realität. Sie sind Ebenen der Validierung durch Verwendung und zeigen die
Ingenieurkonsequenz in konkreten Kontexten.

\section{In den Anwendungen erkennbare Stärken}

Die Kapitel zeigen eine konsistente Anwendbarkeit auf Rekonstruktion,
Auswertung, realisierungsbewusste Interpretation,
Fahrzeugraumanwendungen und Interoperabilitätsszenarien, während sie die
begrifflichen Grenzen wahren.

\section{Grenzen und praktische Bedingungen}

Die Anwendungsqualität hängt von Datenqualität, Kontextverfügbarkeit,
Operatorkonfiguration und Disziplin im Arbeitsablauf ab. Dies sind
Einsatzbedingungen und keine Widersprüche der Begriffsarchitektur.

Offene Fragen der prozesszentrierten Modellierung (einschließlich
PIM-bezogener Fragen) bleiben externe Forschungsthemen und werden in
dieser Migration nicht zur kanonischen Architektur erhoben.

\section{Interoperabilitätsfolge}

BIM/GIS- und Austauschökosysteme werden am besten als nachgelagerte
Repräsentations- und Interoperabilitätsebenen behandelt. So bleibt die
kanonische Zuständigkeit für Ingenieursemantik gewahrt, während die
Integration in Projektarbeitsabläufe möglich wird.

\section{Verbindung zu AIM}
\begin{summary}
Der Anwendungsteil bestätigt, dass Alignment-zentrierte AIM-Begriffe in
vielfältigen Ingenieurkontexten praktisch einsetzbar sind, ohne neue
konkurrierende Begriffsheimaten zu erfordern.

Seine Diskussion betont daher Konsequenz und Grenzdisziplin, anstatt die
Grundlagenforschung erneut zu eröffnen.
\end{summary}
